\section{Spectral decomposition and exact capacity}
\label{gad:section}\label{gcp:section}
\begin{lemma}[Finite approximate incidence decomposition]\label{gad:decomposition}
Let $G_n$ be properly $d_n$-coloured, loopless simple $d_n$-in/out-regular digraphs, with $d_n\ge\alpha{}n$. After taking a subsequence, there are a fixed $k$, a physical set $U$ of size $n-o(n)$, partitions $U=\bigsqcup_iA_i=\bigsqcup_iB_i$, and $o(n)$ exceptional colours such that, with $D=d_n-o(n)$ retained colours and $p=\lim{}d_n/n$:
\begin{enumerate}[label=(\roman*)]
\item $|A_i|,|B_i|=a_in+o(n)$, where $a_i\ge{}p$ and $\sum_ia_i=1$;
\item the retained arcs $A_i\to{}B_i$ give every vertex in/outdegree $D-o(n)$, and each retained colour has $a_in-o(n)$ arcs in block $i$;
\item writing $m_i=\max(|A_i|,|B_i|)$ and padding the real block matrix $Q_i$ by zero rows or columns, $\|Q_i-(D/m_i)J\|_2\le\rho{}D$ for a fixed $\rho<1$;
\item $|B_i\cap{}A_j|/n\to{}c_{ij}$, where $C=(c_{ij})$ is a nonnegative circulation with both margins $(a_i)$.
\end{enumerate}
Thus $k\le1/p$. Deleting the zero-limit physical cells preserves these conclusions, with a possibly larger fixed $\rho<1$.
\end{lemma}
\begin{proof}
Take a subsequence with $d_n/n\to{}p\ge\alpha$. Maintain finitely many real incidence blocks, each with all degrees $d_n-o(n)$ and both sides of size at least $d_n-o(n)$, and $o(n^2)$ total off-block arcs. Initially there is one block. After a cut of size $o(n^2)$, delete both copies of the $o(n)$ physical vertices whose cut degree is more than $o(n)$ (choose the vanishing threshold slowly). All maintained properties survive.

For a current block $Q$ on $X,Y$, counting edges gives $|X|-|Y|=o(n)$. Put $m=\max(|X|,|Y|)$ and $N=R_X^{-1/2}QR_Y^{-1/2}$, where the real degree matrices satisfy $R_X,R_Y=d_nI+o(n)I$ in operator norm. Hence $N=Q/d_n+o(1)$. Its top singular vectors are the normalized square-root degree vectors, whose outer product is $J_{X,Y}/\sqrt{|X||Y|}+o(1)$. Padding the smaller side and replacing this mean by $J_m/m$ costs $o(1)$, so $\sigma_2(N)=\|Q^{\rm{}pad}-(d_n/m)J_m\|_2/d_n+o(1)$ \textup{(AD.1)}. Pass to a further subsequence: either this is bounded below one or it tends to one. Keep the block in the first case. In the second, the real bipartite graph has a normalized-Laplacian eigenfunction $f$ of eigenvalue $\xi=o(1)$ with weighted mean zero and $\sum_{uv\in{}E}(f(u)-f(v))^2=\xi\sum_vd(v)f(v)^2$. Subtract a weighted median; its weighted squared norm cannot decrease. The positive and negative parts have supports of volume at most half the total, and their energies sum to at most the original energy. One nonzero part $g$ has Rayleigh quotient at most $\xi$. For $S_t=\{v:g(v)^2>t\}$, coarea and Cauchy--Schwarz yield $\int{}e(S_t,\overline{S}_t)dt\le\Bigl(\sum_{uv}(g(u)-g(v))^2\sum_{uv}(g(u)+g(v))^2\Bigr)^{1/2}\le\sqrt{2\xi}\int\operatorname{vol}(S_t)dt$. Some nonempty $S=X_1^+\cup{}Y_1^-$ therefore has at most half the volume and boundary $b=o(d_n|S|)$. If $u=|S|$, simplicity gives $b\ge(d_n-o(n))u-u^2/2$, whence $u\ge2d_n-o(n)$. Its complement has linear order and boundary $o(n^2)$; the same inequality gives order at least $2d_n-o(n)$. The two degree sums in $S$ differ by at most $b$, so $|X_1|-|Y_1|=o(n)$ and both sizes are at least $d_n-o(n)$; the complementary sides obey the same bound. Split into these two blocks and perform the physical deletion above.

Each split increases the block count; each block uses at least $d_n-o(n)$ source vertices. There are at most $1/p$ blocks in the limit, so the process stops after boundedly many splits, with every final gap bounded below one. Deleting $o(n)$ rows or columns changes a simple real matrix by Frobenius norm $o(n)$ and its constant term by $o(n)$, so all earlier gaps survive. A common enlarged $\rho<1$ works.

Choose $\varepsilon_n\to0$ slowly enough that the final off-block population is $o(\varepsilon_nn^2)$. Declare a colour exceptional if it has more than $\varepsilon_nn$ off-block arcs; only $o(n)$ colours are exceptional. Each remaining original perfect matching has $|A_i|-o(n)=|B_i|-o(n)$ edges in every block. Physical deletions cost $o(n)$ further edges per colour. Exceptional colours have maximum total row and column degree $o(n)$, so removing them changes each matrix by operator norm $o(n)$. Replacing $d_n$ by $D$ has the same cost. Finally take simultaneous limits of the finitely many side and cell proportions. The two partitions cover the same physical set, giving equal margins and total mass one. Zero-limit cells contain only $o(n)$ vertices and are covered by the same stability estimates.
\end{proof}

\begin{lemma}[An incoming bridge]\label{gcp:bridge}
Fix $\alpha>0$ and $s\ge2$. Suppose $G$ is strongly connected, properly $d$-coloured and $d$-in/out-regular, with $d\ge\alpha{}n$, and $V(G)=U_1\sqcup\cdots\sqcup{}U_s\sqcup{}E$, where $|E|=o(n)$ and every vertex of $U_i$ has internal in/outdegree $d-o(n)$. Every prescribed core has a rainbow simple incoming path of at most three arcs from another core.
\end{lemma}
\begin{proof}
Fix $0<\eta<\alpha/(4s)$ and set $\operatorname{In}(v)=\{i:d^-_{U_i}(v)\ge\eta{}n\}$ and $\operatorname{Out}(v)=\{i:d^+_{U_i}(v)\ge\eta{}n\}$. These sets are nonempty, and both equal $\{i\}$ for $v\in{}U_i$. On $[s]$ put $i\to{}j$, $i\ne{}j$, if some $v$ has $i\in\operatorname{In}(v),j\in\operatorname{Out}(v)$, or some arc $a\to{}b$ has $i\in\operatorname{In}(a),j\in\operatorname{Out}(b)$. In the first case $v\in{}E$ and the $\eta{}n$ choices at both ends give a rainbow two-arc path. In the second choose $w\in{}U_i$ entering $a$ and $u\in{}U_j$ reached from $b$, avoiding the middle colour and each other's colour. Types give $a\notin{}U_j,b\notin{}U_i$; looplessness and these exclusions make $w,a,b,u$ distinct.

This auxiliary digraph is strongly connected. Otherwise choose nonempty proper $I$ with no outgoing edge, put $J=[s]\setminus{}I$, and let $S=\{v:\operatorname{In}(v)\cap{}I\ne\varnothing\}$, $T=\{v:\operatorname{Out}(v)\cap{}J\ne\varnothing\}$. The two edge rules give $S\cap{}T=\varnothing$ and no $S\to{}T$ arc. All cores lie in the appropriate one of $S,T$, so $R=V\setminus(S\cup{}T)\subseteq{}E$. Moreover $e(S,R)\le(s\eta{}n+|E|)|R|,e(R,S)\ge(d-s\eta{}n-|E|)|R|$. Cut balance gives $e(S,R)=e(T,S)+e(R,S)\ge{}e(R,S)$. Since $d>2s\eta{}n+2|E|$, this forces $R=\varnothing$, contradicting strong connectivity of $G$. Every auxiliary vertex thus has an incoming edge; its realization is the required bridge.
\end{proof}

\begin{lemma}[Exceptional colours and exact capacity]\label{gcp:prefix}
Under the core hypotheses above (also allowing $s=1$), let $\mathcal{E}$ be $o(n)$ exceptional colours. Suppose the good-coloured graph inside every core has positive linear families of rainbow connectors between each distinct ordered pair, of bounded endpoint-compatible lengths, tolerating $o(n)$ forbidden vertices and colours. Fix a final core $U$. If $|U|\ge{}d$ or $s\ge2$, there is a rainbow prefix $P$ of $a=o(n)$ arcs ending at $y\in{}U$, using every exceptional colour, such that $d-a+|(V(P)\cap{}U)\setminus\{y\}|\le|U|$ \textup{(GC.1)}. If $|U|\ge{}d$, its end $y$ can be prescribed.
\end{lemma}
\begin{proof}
Let $q=\max(0,d-|U|)=o(n)$. If $q=0$, start with the one-vertex path $y$. Otherwise choose a bridge from another core $W$ into $U$ first, then prepend $q$ fresh good arcs inside $W$. Greedy backward extension is possible since each current vertex has $d-o(n)$ good predecessors and only $o(n)$ vertices and colours are forbidden. The path contains at least $q$ vertices outside $U$.

Suppose its current start is $x$ in a core $I$, and exceptional colour $c$ is unused. The full original colour-$c$ matching has $|I|=\Omega(n)$ arcs entering $I$. At most $|E|$ start outside the cores and at most twice the current path order meet the path. Choose a fresh arc $w\to{}z$ with $w$ in a core and $z\in{}I$, then a good connector $z\to{}x$ inside $I$ avoiding $w$, the previous path except $x$, and previous colours. Its good palette avoids all exceptional colours. Prepending keeps the path simple and rainbow and its start in a core. Process all exceptional colours in this way, costing bounded length each. At the end put $Q=|V(P)\setminus{}U|\ge{}q$. Since $P$ has $a+1$ vertices, $|(V(P)\cap{}U)\setminus\{y\}|=a-Q,d-a+(a-Q)=d-Q\le{}d-q\le|U|$.
\end{proof}

\begin{proposition}[Completion from physical cores]\label{gcp:completion}
Consider a strongly connected counterexample sequence with $d/n\to{}p>0$ and the preceding finite cores, $|U_i|/n\to\beta_i>0$. Suppose all but $o(n)$ colours have $|U_i|-o(n)$ arcs inside every core, and the connector hypotheses of Lemma~\ref{gcp:prefix} hold. Assume that whenever $\beta_i>p$, after deleting any $o(n)$ physical vertices other than a prescribed root and $o(n)$ whole colours, the retained graph in $U_i$ has a rainbow path from that root with one fewer arc than its retained colour count. Then $G$ has a rainbow path of $d-1$ arcs.
\end{proposition}
\begin{proof}
Internal degree gives $\beta_i\ge{}p$. If some $\beta_i>p$, choose that $U_i$; then $|U_i|\ge{}d$. Lemma~\ref{gcp:prefix} gives a prefix of $a=o(n)$ arcs ending at $y$. Delete its earlier vertices and its colours, leaving $D=d-a$ good colours. The assumed rooted suffix has $D-1$ arcs, meets the prefix only at $y$, and has disjoint palette; the concatenation has $d-1$ arcs.

Otherwise all $\beta_i=p$. If $s=1$, then $p=1$ and Theorem~\ref{pf:theorem} applies directly to the original graph. If $s\ge2$, strengthen the prefix capacity by one. When $|U|>d$ this holds by integrality. When $|U|<d$, prepending $q=d-|U|$ arcs before the bridge actually gives $Q\ge{}q+1$, counting also the bridge's initial vertex in the other core. When $|U|=d$, start with an incoming bridge, with no preparatory arcs, then process exceptional colours; here $Q\ge1$. Thus always $|U|-d+Q\ge1$. Set $F=(V(P)\cap{}U)\setminus\{y\}$, $m=|U\setminus{}F|$, and $D=d-a$. Exactly $m-D=|U|-d+Q\ge1,D/m\longrightarrow1$. The retained good graph has in/outdegrees $D-o(n)$ and each colour has $m-o(n)$ arcs. The robust rooted above-half Theorem~\ref{pf:theorem} supplies a suffix of $D-1$ arcs from $y$. Concatenation completes the path.
\end{proof}
